\section{Normalisateur \(899\), unité de Pell et structure arithmétique du mode \(30\)}

Le normalisateur possède la factorisation arithmétique

\begin{equation}
899=30^2-1=29\cdot31.
\label{mab:rm:eq:899}
\end{equation}

Le même mode \(30\) qui engendre les secteurs \(90\) et \(120\) détermine
simultanément la paire de nombres premiers jumeaux \(29,31\). De plus,

\begin{equation}
\varepsilon_{30}=30+\sqrt{899},\qquad
\varepsilon_{30}^{-1}=30-\sqrt{899},\qquad
\log\varepsilon_{30}=\arcosh30.
\label{mab:rm:eq:pell30}
\end{equation}

Cette unité de Pell détermine l’échelle hyperbolique associée au
normalisateur et relie la phase dodécaphasique, le mode \(30\), les nombres premiers
jumeaux et la géométrie hyperbolique. Cette relation n’identifie pas
l’opérateur constitutif du vide au renormalisateur de doublement ; une telle
identification exigerait un entrelaceur explicite.

L’unité agit de manière entière par
\begin{equation}
M_{30}=\begin{pmatrix}30&899\\1&30\end{pmatrix},
\qquad \det M_{30}=1.
\label{mab:rm:eq:pell-matrix30}
\end{equation}
Sur \(\Z[\sqrt{899}]\), cette matrice représente la multiplication par
\(30+\sqrt{899}\). Ses valeurs propres sont \(\varepsilon_{30}^{\pm1}\) ; par
conséquent, \(\log\varepsilon_{30}\) est la longueur de translation de la transformation
hyperbolique associée. Le même entier central \(30\) réapparaît dans le vide
parce que \(90=3\cdot30\) et \(120=4\cdot30\), mais les deux opérateurs continuent
d’avoir des domaines distincts. La confluence démontrée est celle du mode et de ses
complétions ; l’identité des dynamiques n’est pas postulée.

